\section{Threats to Validity}
\label{sec:threats}

The development and synthetic datasets are not statistically representative of
all enterprise environments and cannot by themselves establish broad field
accuracy. This threat is addressed by preserving C1 and C2 as held-out
confirmatory datasets under a pre-registered single-pass protocol.

Construct validity depends on whether operational evidence is mapped to CVSS
v4.0 semantics correctly. The architecture therefore separates control
presence, control effectiveness, and the final Modified metric; it also keeps
organizational Security Requirements distinct from technical resilience
controls. An auditable deterministic deduction can still be wrong if the
underlying rule is wrong.

Internal validity is strengthened by a frozen baseline, versioned rule catalog,
clean-room execution, cryptographic hashes, explicit provenance, and
fail-closed abstention. The confirmatory protocol prohibits result-driven rule
changes, selective retries, or threshold adjustment after exposure.

External validity remains bounded by the diversity of the held-out scenarios
and by the operational sources represented in the evidence model. Results must
not be generalized beyond the tested baseline and datasets.

Conclusion validity is protected by pre-registering the primary C1/C2
acceptance thresholds before unsealing. Confidence intervals and descriptive
statistics are reported as uncertainty summaries and do not replace or alter
the pre-registered pass/fail criteria.

HCAL is an optional secondary layer. Any post-hoc HCAL analysis is reported
separately and cannot change the deterministic confirmatory verdict.
